% Generated draft: requires mathematical and source review.
\documentclass[11pt]{article}
\usepackage[margin=25mm]{geometry}
\usepackage{fontspec}
\setmainfont{Latin Modern Roman}
\usepackage{amsmath,amssymb,mathtools,unicode-math}
\setmathfont{Latin Modern Math}
\usepackage{enumitem,needspace,xurl,hyperref}
\setlength{\parindent}{0pt}
\setlength{\parskip}{5pt}
\setlength{\emergencystretch}{3em}
\newcommand{\N}{\mathbb{N}}
\newcommand{\Z}{\mathbb{Z}}
\newcommand{\Q}{\mathbb{Q}}
\newcommand{\R}{\mathbb{R}}
\newcommand{\C}{\mathbb{C}}
\newcommand{\F}{\mathbb{F}}
\newcommand{\A}{\mathbb{A}}
\newcommand{\PP}{\mathbb{P}}
\newcommand{\E}{\mathbb{E}}
\newcommand{\eps}{\varepsilon}
\newcommand{\dd}{\,\mathrm{d}}
\begin{document}
\section*{UnsolvedMath ID: 30002485}
\textbf{Source problem code:} catalogue entry
\subsection{Definitions and mathematical statement}
The supplied formulation omits the ambient reductive group and the precise homology convention. The following describes the usual interpretation and makes these unresolved qualifications explicit.

Let \(G\) be an ambient real reductive group, let \(\mathfrak g=\operatorname{Lie}(G)\), and let \(K\) be a maximal compact subgroup. The precise class of groups intended by the source is unspecified. The infinitesimal real-sphericity condition on a real Lie subalgebra \(\mathfrak h\subseteq\mathfrak g\) is that there exists a minimal parabolic subalgebra \(\mathfrak q\subseteq\mathfrak g\) such that
\[
\mathfrak g=\mathfrak h+\mathfrak q
\]
as real vector spaces. Whether the source additionally requires \(\mathfrak h\) to arise from a closed or algebraic subgroup is not stated.

Write \(\mathfrak g_{\C}=\mathfrak g\otimes_{\R}\C\) and \(\mathfrak h_{\C}=\mathfrak h\otimes_{\R}\C\). In the usual complex-coefficient interpretation, a Harish--Chandra module \(V\) is an admissible, finitely generated \((\mathfrak g_{\C},K)\)-module: its compatible \(K\)-action is locally finite, and every irreducible \(K\)-type has finite multiplicity. Let \(i:V\hookrightarrow V^{\infty}\) be the inclusion into the smooth moderate-growth completion named in the source. Conventionally this means a smooth Fr\'echet globalization whose \(K\)-finite vectors identify with \(V\); the exact globalization category is not specified in the record.

For clarity, under the ordinary algebraic Lie algebra homology convention, set
\[
C_p(\mathfrak h_{\C},W)=\bigwedge^p\mathfrak h_{\C}\otimes_{\C}W,
\qquad
H_p^{\mathrm{alg}}(\mathfrak h_{\C},W)
=\frac{\ker(d_p)}{\operatorname{im}(d_{p+1})},
\]
where \(d_p\) is the Chevalley--Eilenberg differential for the restricted \(\mathfrak h_{\C}\)-action on \(W\). This quotient uses the actual image, without replacing it by its closure. The inclusion \(i\) induces the chain map \(\operatorname{id}\otimes i\) and hence the comparison map.

\textbf{Comparison question, conditional on these conventions.} For every ambient datum satisfying the intended source hypotheses, every real spherical \(\mathfrak h\), every Harish--Chandra module \(V\), and every integer \(p\geq 0\), is
\[
\Phi_p=H_p^{\mathrm{alg}}(\mathfrak h_{\C},i):
H_p^{\mathrm{alg}}(\mathfrak h_{\C},V)
\longrightarrow
H_p^{\mathrm{alg}}(\mathfrak h_{\C},V^{\infty})
\]
an isomorphism of complex vector spaces?

The source writes this map as \(H_p(\mathfrak h,V)\to H_p(\mathfrak h,V^{\infty})\), without specifying complexification, algebraic versus topological homology, reduced versus unreduced homology, or whether a topological isomorphism is intended. Consequently the displayed interpretation is a conditional statement-only draft, not a resolution of those missing conventions.
\subsection{Short English statement}
The question asks whether passing from a Harish--Chandra module to its smooth moderate-growth completion preserves every nonnegative-degree Lie algebra homology group for a real spherical subalgebra. The ambient group, globalization category, and homology conventions need clarification. The supplied dated curation labels the problem partially solved; that assessment has not been independently verified.
\subsection{Sources}
\begin{itemize}
\item Local supplied UnsolvedMath registry record: permanent ID \(30002485\), title \emph{Homology Comparison for Real Spherical Subalgebras}, exact source problem code \texttt{catalogue entry}, in the Oberwolfach Reports Open Problems collection. The supplied proposed year is 2014; no individual proposer is identified.
\item Supplied source citation: \emph{Representation Theory and Analysis of Reductive Groups: Spherical Spaces and Hecke Algebras} (2014), Oberwolfach Reports, pp.\ 145--212, DOI 10.4171/owr/2014/03. \url{https://doi.org/10.4171/owr/2014/03}. The background describes Oberwolfach Reports as documenting mathematical workshops and their open-problem discussions; no problem-specific page is supplied.
\item Supplied literature reference: Aizenbud--Gourevitch--Kroetz--Liu, homology comparison. \url{https://citeseerx.ist.psu.edu/document?doi=332f58afd544ba575d7c9f67e9d1ec8034b6d52c&repid=rep1&type=pdf}.
\item Supplied literature reference: Li--Liu--Yu, Casselman comparison theorem. \url{https://arxiv.org/abs/2102.03204}.
\item Supplied literature reference: Oberwolfach representation theory report. \url{https://ems.press/content/serial-article-files/46494?nt=1}.
\item Supplied literature assessment, checked on \texttt{2026-08-21}: the comparison theorem is described as proved for standard minimal parabolic nilradicals, while the broader conjecture for real spherical subalgebras is described as not fully established. The supplied status is \texttt{partially\_solved}. This is supplied dated curation derived from the OWR record, not an independently verified literature assessment; the record explicitly calls for independent verification of time-sensitive status and cited claims.
\end{itemize}
\end{document}
